\documentclass[10pt,a4paper]{amsart}
\usepackage[margin=19mm]{geometry}
\usepackage{amsmath,amssymb,mathtools}
\usepackage[scr=rsfs,cal=euler]{mathalfa}
\usepackage{xfrac}
\usepackage[unicode,hidelinks,bookmarks=false]{hyperref}
\newcommand{\ie}{i.e.\ }
\newcommand{\cf}{cf.\ }
\newcommand{\resp}{respectively\ }
\newcommand{\Subsection}{\subsection}
\providecommand{\nobreakdash}{\nobreak\hbox{-}}
\setlength{\emergencystretch}{2em}
\multlinegap0pt
\usepackage{amssymb}
\newtheorem{thm}{Theorem}
\newcommand{\A}{{\mathcal{A}}}
\newcommand{\B}{{\mathcal{B}}}
\renewcommand{\bar}[1]{\overline{#1}}
\DeclareMathOperator*{\bigdoublevee}{\bigvee\mkern-15mu\bigvee}
\renewcommand{\phi}{\varphi}
\providecommand{\tightlist}{%
  \setlength{\itemsep}{0 pt}\setlength{\parskip}{0 pt}}
\title{Relations enumerable from positive information}
\author{Barbara F. Csima, Luke MacLean, Dino Rossegger}
\date{Source: arXiv:2206.01135v3 (CC BY 4.0)}
\begin{document}
\maketitle

\noindent\emph{Source and licence.} Relations enumerable from positive information. Version arXiv:2206.01135v3 (CC BY 4.0), distributed by arXiv under the Creative Commons Attribution 4.0 International licence, http://creativecommons.org/licenses/by/4.0/, as recorded in the arXiv licence metadata for this exact version and shown on the abstract page https://arxiv.org/abs/2206.01135v3. Full licence notice: \texttt{licenses/arxiv-2206.01135-CC-BY-4.0.txt}.\par\smallskip

\noindent\textit{Editorial scope.} Counted unit: the article's thm ripe (source0.tex lines 231--237). The article's complete original proof of it is delivered with it (source0.tex lines 238--276, one \texttt{proof} environment of 39 lines). No mathematics is written, repaired or re-derived: the statement and the proof are the article's own lines, in the article's own notation, and no retained line is substituted or re-flowed.  No further source-local context is needed: every cross-reference of every retained line resolves inside the two delivered spans.  Nothing was omitted from any retained span, so the package carries no omission record.  No retained line contains a citation, so the wrapper carries no bibliography and no bibliography entry.  No retained line contains a figure environment, an image-inclusion command or drawing code, so no image is delivered and the package declares no asset dependency.  Editorial formatting only, in addition: an AMS \texttt{amsart} preamble that restates the article's own declarations and macros used by the retained lines, namely the environments \texttt{thm} on separate counters, together with the standard packages those lines need.  The retained environments print in this document as Theorem 1 in order of appearance, whereas the article prints the same items with the numbering of its own full document.  The complete native source of the article as distributed in the arXiv e-print for this exact version is bundled unchanged as \texttt{sources/126455/source0.tex}.
\begin{thm}\label{ripe}
    Let $\A$ be a structure and $R \subseteq A^{< \omega}$ a relation on it. Then the following are equivalent:
    \begin{enumerate}\tightlist
        \item[(i)] $R$ is relatively intrinsically positively enumerable in $\A$,
        \item[(ii)] $R$ is $\Sigma^p_1$-definable in $\A$ with parameters.
    \end{enumerate}
\end{thm}

\begin{proof} Assuming (ii), there is a uniformly computable sequence $(\phi_i(\overline x,\overline z))_{i\in\omega}$ of $\Sigma_1^p$ formulas where each $\phi_i$ is of the form $\bigdoublevee_{j \in \omega} \exists \overline{y}_{i,j} \psi_{i,j} (\overline{x},\overline{y}_{i,j},\overline{z})$ with the property that for every $\B\cong \A$ there is a tuple $\bar c\in \omega^{|\bar z|}$ such that for all $i\in\omega$ and $\bar a\in\omega^{i}$, $\B\models \varphi_i(\bar a,\bar c)$ if and only if $\bar a \in R^{\B}$.
%\begin{equation*}
%    R(\overline{x}) \Leftrightarrow \varphi(\overline{x},\overline{z}) \equiv \bigdoublevee_{i \in I} \exists \overline{y}_i \varphi_i (\overline{x},\overline{y}_i,\overline{z})
%\end{equation*}
%where each $\varphi_i$ is quantifier free and $\overline{c}$ will become our parameterizing tuple.
Recall that each $\psi_{i,j}$ is a finite conjunction of atomic formulas, i.e., $\psi_{i,j}=\theta_1(\overline x, \overline y_{i,j},\overline z)\land \dots \land \theta_n(\overline x, \overline y_{i,j},\overline z)$ for some $n\in \omega$. For $\theta(\overline x)$ an atomic formula, let $\ulcorner \theta(\overline a)\urcorner$ be the function mapping $\theta(\overline a)$ to its code in the positive diagram of a structure. For example, if $\theta(\overline x)=R_i(x_3,x_5)$, then $\ulcorner \theta(\overline a)\urcorner=\langle i+2, \langle a_3,a_5\rangle\rangle$ for $\overline a\in \omega^{<\omega}$. Consider the set $X^{\overline a, \overline b,\overline c}_{i,j}=\{\ulcorner \theta_k(\overline a,\overline b,\overline c)\urcorner: k<n\}$. Clearly, $X^{\overline a,\overline b,\overline c}_{i,j}\subseteq P(\A)$ if and only if $\A\models \psi_{i,j}(\overline a,\overline b,\overline c)$ for any $L$-structure $\A$. We define an enumeration operator $\Psi$ by enumerating all pairs %$(\overline{a},\overline{b})$, and anytime we see $\psi_{i,j}( \overline{a}, \overline{b}, \overline{c})$ being satisfied for one of the formulas $\psi_{i,j}$, some tuples $\overline{a},\overline{b}$, and our fixed tuple $\overline{c}$, we enumerate the pair $ \langle \overline{a}, X_{i,j}^{\overline{a}, \overline{b}, \overline{c}} \rangle$ into $\Psi$. %
$\langle \overline a,X_{i,j}^{\overline a,\overline b,\overline c}\rangle$ into $\Psi$. It follows that 
\begin{align*}
    \overline a\in \Psi^{P(\B)}&\Leftrightarrow \exists \langle \overline a,X_{|\overline a|,j}^{\overline a, \overline b, \overline c}\rangle \in \Psi \land X_{|\overline a|,j}^{\overline a, \overline b\overline c}\subseteq P(\B)\\
    &\Leftrightarrow \B\models \exists \overline y_{|\overline a|,j}\psi_{|\overline a|,j}(\overline a, \overline y_{|\overline a|,j},\overline c)\\
    &\Leftrightarrow \B\models \phi_{|\overline a|}(\bar a, \bar c)
\end{align*}
and thus $R$ is r.i.p.e.
%Thus $\overline{a}$ is in $R^\B$, if and only if there is some $j$ and $\overline{b}$ such that $\B \models \psi_{|\overline a|,j}(\overline{a},\overline{b},\overline{c})$ if and only if the pair $ \langle \overline{a}, X_{|\overline a|,j}^{\overline{a}, \overline{b}, \overline{c}} \rangle$ is in $\Psi$ if and only if $\overline{a} \in \Psi^{P(\B)}$.

%To show that (i) implies (ii) we build an enumeration $g:\omega \rightarrow \A$ by constructing a nested sequence of tuples $\{\overline{p}_s\}_{s \in \omega} \subseteq A^{< \omega}$ and letting $g = \bigcup_s \overline{p}_s$. We then define $\B=g^{-1}(\A)$.
Let $g$ be a r.i.p.e.\ generic enumeration of $\A$ produced as in the paragraph above the theorem and let $\B=g^{-1}(\A)$. Our goal is to produce a $\Sigma^p_1$ definition of $R^\B$, as $\B\cong \A$ we get that this is then also a definition for $R$. %To obtain it we try to force during the construction that $R^\B=g^{-1}(R)\neq \Psi_e^{P(\B)}$. As $R$ is r.i.p.e.\ this will fail for some $e$ and we will use this failure to get a syntactic definition.
%
%To construct $\B$ we do the following. Let $\overline{p}_0$ be the empty sequence. At stage $s+1=2e+1$ if the $e$-th element of $A$ is not already in $\overline{p}_s$ then we let $\overline{p}_{s+1}=\overline{p}_s^\frown e$. Otherwise let $\overline{p}_{s+1}=\overline{p}_s$. This guarantees that $g$ is onto.
%At stage $s+1=2e$ we try to force $\Psi_e^{P(\B)}\not \subseteq g^{-1}(R)$ for which we need a tuple $\langle j_1,\dots, j_l \rangle \in \Psi_e^{P(\B)}$ with $\langle g(j_1),\dots, g(j_l) \rangle \not \in R$. To do this we ask if there is an extension $\overline{q}$ of $\overline{p}_s$ in the set
Towards this, consider the set
\begin{equation*}
    Q_e = \{ \overline{q} \in A^* : \exists l, j_1 ,\dots j_l < |\overline{q}| \, [\langle j_1,\dots, j_l \rangle \in \Psi_e^{P_\A(\overline{q})} \text{ and } \langle q_{j_1},\dots, q_{j_l} \rangle \not \in R] \}.
\end{equation*}
We have that $g$ meets $Q_e$ if and only if $R^\B\neq \Psi_e^{P_\B}$. By our assumption that $R$ is r.i.p.e., there is $e_0$ such that $\Psi^{P(\B)}_{e_0}=R^\B$, and thus $g$ does not meet $Q_{e_0}$.

%If at any stage $s+1 = 2e$ we succeed in defining $\overline{p}_{s+1}=\overline{q}$ for some $\overline{q} \in Q_e$ then we will have made $\Psi_e^{P(\B)} \neq g^{-1}(R)$. But by our assumption this must fail for some $e \in \omega$. For this $e$, at stage $s+1=2e$ we will not have been able to find an extension of $\bar p_{s}$ in $Q_e$. 
We will use this to give a $\Sigma^p_1$ definition of $R$ with parameters $\bar p_{s}$, where $\bar p_{s}$ is the stage in the construction of $g$ at which we decide $Q_{e_0}$.

Notice that if there is some $\bar q \supseteq \bar p_s$ and sub-tuple $\langle q_{j_1} , \dots, q_{j_l} \rangle$ such that $\langle j_1,\dots , j_l \rangle \in \Psi_e^{P_\A(\overline{q})}$ then we must have $\langle q_{j_1} , \dots, q_{j_l} \rangle \in R$ or else we will have that $\bar q \in Q_{e_0}$. We now show that $R$ is equal to the set
\begin{align*}
    S = \{ \langle q_{j_1},\dots , q_{j_l} \rangle \in A^* : \text{ for some } \bar q \in A^{< \omega} \text{ and } l , j_1, \dots , j_l < |\bar q | \text{ satisfying } \bar q \supseteq \bar p_s \\ \text{ and } \langle j_1 ,\dots , j_l \rangle \in \Psi_e^{P_\A(\bar q)}\}
\end{align*}
By the previous paragraph, $S \subseteq R$. If $\bar a \in R$ then there are indices $j_1,\dots, j_{|\bar a|}$ such that $\bar a = \langle g(j_1), \dots g(j_{|\bar a|}) \rangle$ and so if we take a long enough initial segment of $g$ it will witness the fact that $\bar a \in S$. Fix an enumeration $(\phi_i^{at})_{i\in\omega}$ of all atomic formulas where without loss of generality the free variables of $\phi_i^{at}$ are a subset of $\{x_0,\dots, x_{i}\}$. The following is a $\Sigma^p_1$ definition of $S$
\[
    \bigdoublevee_{ C \subset_{\text{fin}} \omega} \bigvee_{\langle j_1,\dots , j_{|\bar a|} \rangle \in W_e^C} \exists\, \bar q \supseteq \bar p_s \, [\langle q_{j_1}, \dots ,q_{j_{|\bar a|}} \rangle = \bar a \, \wedge \, \bigwedge_{i \in C} [\varphi^{at}_i]\frac{q_0}{x_0}\cdots\frac{q_{|\bar q|}}{x_{|\bar q|}} ]
\]
where the latter half of the formula is simply saying that $C \subseteq P_\A(\bar q)$.
\end{proof}

\end{document}
